Posem dos cables prims conductors sobre una taula perpendiculars entre si i sense que hi hagi contacte elèctric entre ells. Posteriorment, col·loquem un petit imant, una brúixola, a un metre de la taula just per sobre l'encreuament dels dos fils conductors, com indica la figura.

\figura[0.5]{fig1.png}

\begin{apartat}{1,25}
Representeu els camps magnètics creats pels fils $A$ i $B$ en el punt on està situada la brúixola. Si pel fil $A$ hi circula un corrent d'intensitat 5~A, quina intensitat ha de circular pel fil $B$ perquè la brúixola quedi orientada paral·lela al fil $B$?
\end{apartat}

\begin{apartat}{1,25}
Pel fil $A$ hi passa una intensitat $I_A = 10\ \mathrm{A}$ i la brúixola queda orientada amb un angle de $30^\circ$ respecte al fil $B$. Quina intensitat passa pel fil $B$?
\end{apartat}
